% !TeX program = lualatex
% !TeX encoding = utf8
% !TeX spellcheck = uk_UA
% !TeX root =../PyscfBook.tex

%% --------------------------------------------------------
\section{Основи методу конфігураційної взаємодії}
%% --------------------------------------------------------


Метод \textbf{конфігураційної взаємодії (Configuration Interaction, CI)} --- це \textbf{варіаційний пост-Хартрі–Фоківський підхід}, у якому хвильова функція системи представляється як \textbf{лінійна комбінація детермінантів Слейтера}, що відповідають різним конфігураціям розподілу електронів по орбіталях. Таким чином, ураховується змішування кількох електронних конфігурацій, що дозволяє відтворити ефекти електронної кореляції, відсутні в однодетермінантному наближенні Хартрі–Фока.

\begin{equation}
    |\Psi_{\text{CI}}\rangle =
    c_0 |\Phi_0\rangle +
    \sum_i^{\text{occ}} \sum_a^{\text{virt}} c_i^a |\Phi_i^a\rangle +
    \sum_{i<j} \sum_{a<b} c_{ij}^{ab} |\Phi_{ij}^{ab}\rangle + \ldots
\end{equation}

де:
\begin{itemize}
    \item $|\Phi_0\rangle$ --- детермінант Хартрі–Фока (основна конфігурація);
    \item $|\Phi_i^a\rangle$ --- \textbf{одинично збуджений} (singly excited) детермінант, у якому один електрон переходить із зайнятої орбіталі $i$ на віртуальну орбіталь $a$;
    \item $|\Phi_{ij}^{ab}\rangle$ --- \textbf{подвійно збуджений} (doubly excited) детермінант --- два електрони переходять із орбіталей $i,j$ на $a,b$;
    \item $|\Phi_{ijk}^{abc}\rangle$ --- \textbf{потрійно збуджений} (triply excited) детермінант;
    \item і т. д. --- вищі збудження (quadruple, quintuple, \ldots).
\end{itemize}

\paragraph{Фізичний зміст методу.}
Ідея CI полягає у варіаційному визначенні коефіцієнтів $c_0, c_i^a, c_{ij}^{ab}, \ldots$ шляхом мінімізації середнього значення енергії Гамільтоніана:
\[
E_{\text{CI}} = \frac{\langle \Psi_{\text{CI}} | \hat{H} | \Psi_{\text{CI}} \rangle}{\langle \Psi_{\text{CI}} | \Psi_{\text{CI}} \rangle}.
\]
Отримана хвильова функція є лінійною комбінацією конфігурацій, що найкраще описує істинний багаточастинковий стан системи в межах обраного базисного набору. Чим більше збуджень враховано, тим точніше відтворюється електронна кореляція.

\paragraph{Основні варіанти методу CI:}
\begin{itemize}
    \item \texttt{CIS} (Configuration Interaction Singles) --- \textbf{метод одиночних збуджень}: ураховуються лише $|\Phi_i^a\rangle$. Використовується переважно для опису збуджених електронних станів, але не покращує енергію основного стану.
    \item \texttt{CISD} (Configuration Interaction Singles and Doubles) --- \textbf{одиничні та подвійні збудження} ($S+D$). Найуживаніший практичний варіант, що враховує значну частину динамічної кореляції.
    \item \texttt{CISDT} --- \textbf{одиничні, подвійні та потрійні збудження} ($S+D+T$); точніший, але обчислювально значно дорожчий.
    \item \texttt{Full CI} (повна конфігураційна взаємодія) --- ураховуються \textbf{всі можливі збудження} у межах вибраного базису. Це \emph{точне розв’язання} рівняння Шредінґера для заданого базису, яке слугує еталоном для інших методів.
\end{itemize}

\paragraph{Реалізація методів CI у програмному пакеті \texttt{PySCF}}
У системі \texttt{PySCF} реалізовано кілька основних варіантів конфігураційної взаємодії, які охоплюють найпоширеніші рівні наближення:
\begin{itemize}
    \item \texttt{CISD} (\inlinecode{pyscf.ci.CISD}) --- метод одиночних і подвійних збуджень, що базується на детермінанті Хартрі–Фока. Це найуживаніша реалізація CI у PySCF; підтримуються обмежені (RHF), нерестриковані (UHF) та загальні (GHF) варіанти хвильової функції.
    \item \texttt{Full CI} (\inlinecode{pyscf.fci.FCI}) --- повна конфігураційна взаємодія у межах заданого базису. Використовується для малих систем як еталонне («точне») розв'язання рівняння Шредінґера. Дозволяє задавати спін, симетрію хвильової функції, «заморожувати» орбіталі (frozen core) та обчислювати багаточастинкові щільності.
\end{itemize}
Інші варіанти CI (\texttt{CISDT}, \texttt{CISDTQ} тощо) у \texttt{PySCF} \textbf{не реалізовані безпосередньо}, оскільки їхня обчислювальна складність різко зростає з розміром системи. Для опису більш високих збуджень і точнішого врахування електронної кореляції в \texttt{PySCF} зазвичай використовують методи зв'язаних кластерів (\texttt{CCSD}, \texttt{CCSD(T)}) або багатоконфігураційні підходи (\texttt{CASSCF}, \texttt{MRCI}).

\subparagraph{Приклад використання методів CI}
Нижче наведено типовий приклад розрахунку методами CISD та Full CI для молекули \ce{H2}:

\begin{minted}{python}
from pyscf import gto, scf, ci, fci

# Молекула H2
mol = gto.M(atom='H 0 0 0; H 0 0 0.74', basis='sto-3g')

# Hartree-Fock
mf = scf.RHF(mol).run()
e_hf = mf.e_tot
print(f"E(HF) = {e_hf:.6f} a.u.")

# CISD - метод kernel() повертає (E_corr, civec)
mycisd = ci.CISD(mf)
e_corr_cisd, civec_cisd = mycisd.kernel()
e_tot_cisd = e_hf + e_corr_cisd  # Повна енергія = HF + кореляція
print(f"E(CISD) corr = {e_corr_cisd:.6f} a.u.")
print(f"E(CISD) total = {e_tot_cisd:.6f} a.u.")

# Full CI - kernel() повертає (E_tot, civec)
myfci = fci.FCI(mf)
e_tot_fci, civec_fci = myfci.kernel()
e_corr_fci = e_tot_fci - e_hf
print(f"E(FCI) total = {e_tot_fci:.6f} a.u.")
print(f"E(FCI) corr = {e_corr_fci:.6f} a.u.")
\end{minted}

\textbf{Ключові особливості API:}
\begin{itemize}
    \item \texttt{CISD}: метод \inlinecode{.kernel()} повертає \textbf{кореляційну енергію} та CI-вектор. Повна енергія обчислюється як $E_{\text{tot}} = E_{\text{HF}} + E_{\text{corr}}$.
    \item \texttt{Full CI}: метод \inlinecode{.kernel()} повертає \textbf{повну енергію} та CI-вектор. Кореляційна енергія обчислюється як $E_{\text{corr}} = E_{\text{tot}} - E_{\text{HF}}$.
    \item CI-вектор (\inlinecode{civec}) містить коефіцієнти розкладу хвильової функції за базисом конфігурацій (детермінантів Слейтера). Для простого розрахунку енергії його можна ігнорувати.
\end{itemize}


\paragraph{Недоліки методу.}
Основна проблема CI --- \textbf{комбінаторне зростання числа детермінантів} із розміром системи. Для $N$ електронів і $M$ орбіталей кількість можливих конфігурацій зростає як $\binom{M}{N}$, що робить повний CI практично недосяжним для систем із більш ніж кількома електронами. Тому на практиці використовують обмежені варіанти (CISD, CISDT тощо) або альтернативні методи --- наприклад, \textbf{методи зв’язаних кластерів (Coupled Cluster, CC)}, які зберігають подібну точність за меншої обчислювальної складності.


%\paragraph{Приклад розрахунку методом CISD для атома \ce{He}}
%
%
%\inputcode[mode=inline]{CISD_basic.py}
%
%Розрахунок енергії атома \ce{He} показує, що метод Хартрі–Фока суттєво
%переоцінює енергію системи, оскільки не враховує електронну кореляцію:
%різниця з експериментом становить близько \(\Delta E \approx 49~\text{mHa}\).
%Метод \texttt{CISD} частково враховує кореляційні ефекти,
%зменшуючи похибку до \(\approx 16~\text{mHa}\).
%Однак навіть \texttt{CISD} не досягає межі \texttt{full CI}, що відповідає експериментальному значенню
%\(-2.903724~\text{Ha}\).


%% --------------------------------------------------------
\subsection{Full CI --- точний розв'язок}
%% --------------------------------------------------------

Full CI включає всі можливі детермінанти і дає точну хвильову функцію та енергію в межах обраного базису:

\begin{equation}
    E_{FCI} = \min_{\{c_i\}} \langle \Psi_{FCI} | \hat{H} | \Psi_{FCI} \rangle
\end{equation}

\textbf{Проблема:} Кількість детермінантів зростає як:
\begin{equation}
    N_{det} = \binom{N_{orb}}{N_{\alpha}} \times \binom{N_{orb}}{N_{\beta}}
\end{equation}

Для 10 електронів у 20 орбіталях: $N_{det} \approx 10^{10}$ детермінантів!

Для демонстрації можливостей методу проведемо розрахунок для атома \ce{He}:
\begin{itemize}
    \item \ce{He} --- це найпростіша багатоефектронна система (два електрони), у якій вже проявляється електронна кореляція, відсутня у одноелектронних системах типу \ce{H} чи \ce{H2+};
    \item розмір простору детермінантів для \ce{He} невеликий, тому Full CI можна реалізувати навіть у середніх базисах без великих обчислювальних витрат;
    \item саме на прикладі \ce{He} зручно \textbf{демонструвати вплив базисного набору} на енергію Full CI, оскільки результат швидко сходиться до межі повного базису (\emph{Complete Basis Set Limit}).
\end{itemize}

\inputcode[mode=inline]{FullCI_basic.py}

\begin{commentbox}
Як видно з виводу програми, навіть метод \textbf{Full CI}, який у межах даного базисного набору забезпечує
математично точне розв'язання рівняння Шредінґера, не відтворює експериментальну енергію.
Причина полягає не в обмеженнях самого CI, а у \textbf{неповноті базису}.
У малому базисі (наприклад, STO-3G) атомні орбіталі лише грубо апроксимують
реальні електронні хвильові функції, що призводить до недооцінки електронної кореляції
та деякого завищення енергії (менш від’ємного значення).

Зі збільшенням розміру базису (наприклад, 6-31G, cc-pVTZ, aug-cc-pVQZ)
розрахунок Full CI поступово наближається до \emph{межі повного базису}
(\emph{complete basis set limit}), де відхилення від експерименту стає мінімальним.

Таким чином, \textbf{Full CI є точним лише в межах обраного базису},
але не є абсолютно точним у фізичному сенсі.
\end{commentbox}


%% --------------------------------------------------------
\subsection{CISD розрахунки}
%% --------------------------------------------------------

Метод \texttt{CISD} (\emph{Configuration Interaction with Singles and Doubles})
є наближенням до Full CI, в якому враховуються лише \textbf{одиночні} (S) та \textbf{подвійні} (D) збудження
відносно детермінанта Гартрі–Фока:

\begin{equation}
|\Psi_{\text{CISD}}\rangle = c_0|\Phi_0\rangle
+ \sum_i^N \sum_a c_i^a |\Phi_i^a\rangle
+ \frac{1}{2}\sum_{ij}^{N} \sum_{ab} c_{ij}^{ab} |\Phi_{ij}^{ab}\rangle
\end{equation}
%
де $|\Phi_i^a\rangle$ --- детермінант із одиночним збудженням електрона з орбіталі $i$ у $a$,
а $|\Phi_{ij}^{ab}\rangle$ --- детермінант із подвійним збудженням.

Метод CISD враховує основну частину кореляційної енергії,
зокрема короткодійну кулонівську взаємодію між електронами,
і є \textbf{гарним компромісом} між точністю та обчислювальними витратами.

Проте CISD має фундаментальне обмеження --- \textbf{він не є розмірно узгодженим (size-consistency)}.
Це означає, що для двох незалежних систем A і B
виконується:
\[
E_{\text{CISD}}(\text{A} + \text{B}) \neq E_{\text{CISD}}(\text{A}) + E_{\text{CISD}}(\text{B})
\]
на відміну від методів типу \texttt{MP2} або \texttt{FCI}.

Фізично це пов’язано з тим, що CISD включає лише збудження з одного
детермінанту. Коли система розпадається на незалежні частини, такого
опису недостатньо, щоб відтворити всі можливі комбінації збуджень для
обох фрагментів. Тобто, Метод \texttt{CISD} покращує енергію зв’язку
порівняно з \texttt{HF}, але не відтворює правильну енергію при дисоціації.

Цю властивість можна продемонструвати на прикладі молекули \ce{LiH},
де видно, що на далекій відстані між атомами
енергія \texttt{CISD} є сумою енергій окремих атомів.
Для демонстрації використовують такий тест:

\inputcode[mode=inline]{LiH_size_consistency.py}

Результати демонструють порушення розмірної узгодженості:

\begin{minted}{python}
=================================================================
Метод       | LiH bound | LiH diss  | Li + H   | Error (diss - 2atoms)
=================================================================
HF          | -7.71083  | -7.23965  | -7.78211 |  0.54246
CISD        | -7.79875  | -7.28025  | -7.78242 |  0.50217
Full CI     | -7.79884  | -7.78242  | -7.78242 | -0.00000
Експеримент | -8.07000  | -7.97800  | -7.97800 |  0.00000
=================================================================
\end{minted}

Видно, що при дисоціації молекули \ce{LiH} енергія,
обчислена в межах CISD, не збігається з сумою енергій окремих атомів
\ce{Li} та \ce{H}. У наведеному прикладі похибка становить близько
0.5~Ha, що є суттєвим відхиленням. Для порівняння, повний CI (FCI)
повністю узгоджений із фізичною картиною --- похибка практично нульова.

Цей приклад особливо зручний для демонстрації залежності точності FCI
та наближених методів (як CISD) від базисного набору: зі збільшенням
розмірності базису відхилення CISD зменшується, проте ніколи не
зникає повністю, оскільки сама природа методу не гарантує розмірної
узгодженості.